\documentclass[12pt,a4paper]{extarticle}
\usepackage[utf8]{inputenc}
\usepackage[T5]{fontenc}
\usepackage{geometry}
\geometry{a4paper, left=1.5cm, right=1cm, top=1.5cm, bottom=1.5cm, headheight=20pt, headsep=15pt, footskip=28pt}
\usepackage{enumitem}
\usepackage{amsmath}
\usepackage{amssymb}
\usepackage{diagbox}
\usepackage[most]{tcolorbox}
\newcommand{\khongxacdinh}{\text{KXĐ}}
\newcommand{\blackbox}[1]{\texorpdfstring{\fcolorbox{black}{black}{\textcolor{white}{#1}}}{#1}}
\usepackage{tikz}
\definecolor{bookblue}{RGB}{58,90,172}
\definecolor{book}{RGB}{58,90,172}
\definecolor{myblue}{RGB}{200,40,40}
\newcommand{\bluecheck}{\textcolor{myblue}{\checkmark}}
\newtcolorbox{formulaBox}{colback=gray!5,colframe=bookblue,boxrule=0.8pt,arc=2mm}
\newtcolorbox{notebox}{colback=blue!3,colframe=myblue,boxrule=0.8pt,arc=2mm}
\newtcolorbox{doublelinebox}{
    enhanced,
    colback=blue!2,
    colframe=bookblue,
    boxrule=0.8pt,
    arc=2mm,
    borderline={0.8pt}{2pt}{bookblue}
}
\usetikzlibrary{patterns, patterns.meta} % Thêm patterns.meta vào đây
% Gọi package nằm trong thư mục con template
\usepackage{../template/mngocbox}

\begin{document}

% Sử dụng ngoặc vuông [...] để thay đổi linh hoạt các thông số
\makeheadbox[headerright={Chương 1},chude={Bài 2},tieude={Giá trị lớn nhất và giá trị nhỏ nhất của hàm số},dinhdang={Giải tích},lop={12 },gv={Nguyễn Lê Minh Ngọc}]

\vspace{2em}
\makeheading{I}{Kiến thức cần nhớ}
\subsection*{1. Định nghĩa}
Cho hàm số $f(x)$ xác định trên $D$.
\begin{doublelinebox}
    Số $M$ được gọi là giá trị lớn nhất của hàm số $y = f(x)$ trên $D$, kí hiệu $M = \max_{D} f(x)$ nếu:
    \[
    \begin{cases}
        f(x) \le M, \quad \forall x \in D \\
        \exists x_0 \in D \mid f(x_0) = M
    \end{cases}
    \]
    
    Số $m$ được gọi là giá trị nhỏ nhất của hàm số $y = f(x)$ trên $D$, kí hiệu $m = \min_{D} f(x)$ nếu:
    \[
    \begin{cases}
        f(x) \ge m, \quad \forall x \in D \\
        \exists x_0 \in D \mid f(x_0) = m
    \end{cases}
    \]
\end{doublelinebox}
\subsection*{2. Tìm giá trị lớn nhất và nhỏ nhất của hàm số liên tục}
\noindent $\bullet$ \textbf{Lưu ý:} Mọi hàm số liên tục trên một đoạn đều có giá trị lớn nhất và giá trị nhỏ nhất trên đoạn đó.
\vspace{0.3cm}
\noindent Giả sử hàm số $f(x)$ liên tục trên đoạn $[a; b]$ và có đạo hàm trên khoảng $(a; b)$, có thể trừ đi một số hữu hạn điểm. Nếu $f'(x) = 0$ chỉ tại một số hữu hạn điểm thuộc khoảng $(a; b)$ thì ta có quy tắc tìm giá trị lớn nhất và giá trị nhỏ nhất của hàm số $f(x)$ trên đoạn $[a; b]$ như sau:
\vspace{0.3cm}
\begin{doublelinebox}
    \noindent \textbf{\color{blue}Bước 1.} Tìm các điểm $x_1, x_2, \dots, x_n$ thuộc khoảng $(a; b)$ mà tại đó hàm số có đạo hàm bằng $0$ hoặc không tồn tại đạo hàm.\\
    
    \noindent \textbf{\color{blue}Bước 2.} Tính $f(x_1), f(x_2), \dots, f(x_n), f(a)$ và $f(b)$.\\
    
    \noindent \textbf{\color{blue}Bước 3.} So sánh các giá trị tìm được ở \textit{bước 2}.\\
    
    \noindent Số lớn nhất trong các giá trị đó là giá trị lớn nhất của hàm số $f(x)$ trên đoạn $[a; b]$, số nhỏ nhất trong các giá trị đó là giá trị nhỏ nhất của hàm số $f(x)$ trên đoạn $[a; b]$.
\end{doublelinebox}
\subsection*{3. Lưu ý}

\noindent \textbf{\color{blue}a) Lưu ý về giá trị lớn nhất và nhỏ nhất của hàm số $f(x) = a\sin x + b\cos x$.}\\
Cho hàm số $f(x) = a\sin x + b\cos x$ ($a^2 + b^2 > 0$). Khi đó:
\[
\max f(x) = \sqrt{a^2 + b^2}; \quad \min f(x) = -\sqrt{a^2 + b^2}
\]
Hệ quả: $\max(a\sin x) = |a|$; \quad $\min(a\sin x) = -|a|$.

\vspace{0.4cm}

\noindent \textbf{\color{blue}b) Lưu ý về giá trị lớn nhất và nhỏ nhất của hàm số $f(u(x))$ trên $D$.}\\
Đặt $t = u(x)$, với $x \in D$, giả sử ta tìm được tập giá trị của $u(x)$ trên $D$ là $K$. Khi đó:
\[
\max_{x \in D} f(u(x)) = \max_{t \in K} f(t)
\]

\vspace{0.4cm}

\noindent \textbf{\color{blue}c) Bất đẳng thức Cauchy (BĐT AM-GM) cho 2 số và cho 3 số}\\
Cho $a, b$ là các số thực không âm, khi đó:
\[
\frac{a + b}{2} \ge \sqrt{ab}
\]
Cho $a, b, c$ là các số thực không âm, khi đó:
\[
\frac{a + b + c}{3} \ge \sqrt[3]{abc}
\]
\noindent \textbf{Hệ quả:}
\[
ab \le \left(\frac{a + b}{2}\right)^2, \quad abc \le \left(\frac{a + b + c}{3}\right)^3 \quad \text{với } a, b, c \ge 0
\]

\end{document}